\begin{table}[!htbp]\centering\scriptsize\setlength{\tabcolsep}{3pt}
\caption{Hypothèses fortes et chiffres non vérifiés}\label{tab:ircsg_hyp}
\begin{tabular}{@{}p{0.43\linewidth}p{0.2\linewidth}p{0.33\linewidth}@{}}
\toprule
Hypothèse & Où elle joue & Ce qu'elle change si elle est fausse \\
\midrule
Perception de la CSG : la CSG et la CRDS d'aujourd'hui font réagir le revenu d'activité autant que l'IR. Aucune étude ne le mesure. & Tout le gain sur la marge intensive & Décisif. CSG non perçue : gain $-$2,5~Md€ au lieu de 4,7 ; perçue à moitié : 1,0. \\[2pt]
Élasticités : revenu réel 0,05 (moitié du bas) et 0,10 (au-dessus) ; participation 0,25 sous 2 Smic ; ce sont des choix, pas des estimations françaises. & Tous les gains & Fourchette 2,1 à 10,3~Md€ ; long terme 6,7~Md€. \\[2pt]
Les pensions, le chômage et l'épargne ne réagissent pas (âge de départ, surcote, cumul emploi-retraite, sortie en capital). & Mécanisme du financement par les pensions & Non chiffré ; sens ambigu (une pension nette plus faible pousse à partir plus tard ; un taux marginal élevé dans la zone de sortie de l'exonération pousse en sens inverse). \\[2pt]
Régime permanent : les actifs ne voient pas que leurs futures pensions seront plus taxées. & Part « permanente » du gain & Vu en entier : gain permanent 3,0~Md€ au lieu de 4,7. \\[2pt]
Répartition des revenus au PFU entre tranches (non publiée) ; convention « mixte ». & Charge du 1~\%, forme du haut du barème & Convention « haut » (que soutient un calage sur l'IR observé) : 1~\% $-$3,0~Md€ ; charge inchangée si 50~\% dès 91~k€. \\[2pt]
Foyers purs : pas de foyer mixte (couple dont un seul est retraité, retraité bailleur) ; loyers tous chez les actifs. & Coût de la franchise, charge des retraités & Franchise sous-estimée et hausse des retraités surestimée, jusqu'à 3~Md€ environ (borne grossière). \\[2pt]
Partage des revenus dans les couples (Insee Première 1492, quatre classes ; le même pour les retraités). & Coût de l'impôt individuel, charge des couples de retraités & Couples de retraités à 62/38 : charge des retraités plus faible de 2~Md€ environ (registre, H11, ancienne proposition). \\[2pt]
Enfants : un par parent isolé, deux par couple « avec enfants », aucun chez les retraités et les concubins. & Coût de la franchise par enfant, familles nombreuses & Familles de trois enfants absentes des agrégats (cas-types seulement). \\[2pt]
Charge actuelle par type de foyer répartie selon des cas-types sans niche. & Qui paie entre types de foyers & Charge actuelle des retraités (50,2~Md€) non observable ; niches réparties au prorata. \\[2pt]
Base 2025 : l'IR de 2025 est l'impôt encaissé dans l'année (comptes nationaux), dont la hausse (+7,9~\%) tient en partie au jeu des acomptes et des soldes du prélèvement à la source ; chaque agrégat de 2024 est porté en 2025 avec la même croissance dans toutes les tranches. & Recette à remplacer, charge actuelle par tranche & Si cette hausse est passagère, la recette est surestimée d'environ 4,5~Md€ (IR porté en 2025 comme les revenus), soit environ 0,4 point de taux. \\[2pt]
Épargne retraite : déduction de 10,6~Md€ (cases 2042) supposée proportionnelle aux revenus d'activité de la moitié haute. & Barème, haut de la distribution & Coût de la déduction 3,0~Md€ (0,29~point) ; répartition non vérifiée. \\[2pt]
Cas-types : législation des revenus 2025 d'OpenFisca, Smic de 21\,622~€ (sources secondaires) ; retraités en régime stable (taux de CSG le plus élevé dont le seuil est dépassé par le revenu fiscal calculé à ce taux). & Tous les cas-types, courbes de taux marginaux & Dans une bande de quelques centaines d'euros autour de chaque seuil de CSG, aucun taux n'est cohérent avec son propre revenu fiscal ; le taux inférieur est retenu, ce qui déplace le saut de moins de 1\,000~€. \\[2pt]
Valeur actualisée : table de Gompertz calée sur l'espérance de vie à 65 ans (21,8 ans) ; cohortes de retraités de même taille. & Valeur actualisée de la hausse des retraités & Ordre de grandeur ; 1 à 3~\% réel donne la fourchette. \\[2pt]
Chiffres non vérifiés sur une source primaire : CSG des pensions de 2024 (CCSS) ; croissance de la CSG par assiette en 2025 (régimes de base et FSV, appliquée à tous les affectataires) ; paramètres 2025 d'OpenFisca (seuils de CSG des pensions, abattement des plus de 65 ans) ; Smic 2025 ; part des familles au plafond du quotient familial ; partage des niches par tranche ; l'IRCOM inclut l'IR au PFU et exclut la CEHR. & Calage, cas-types, récit & Registre revue/hypotheses\_ircsg.md, section 4. \\[2pt]
\bottomrule
\end{tabular}
\par\vspace{3pt}\parbox{\linewidth}{\raggedright\scriptsize\color{encre}
Source : registre revue/hypotheses\_ircsg.md (hypothèses H1 à H20), chiffres recalculés pour la proposition par transition\_ircsg.py, controle\_hauts\_purs.py et proposition\_ircsg.py.}
\end{table}
